\section{Empirical Results}

\subsection{Primary Performance and Representation Ablation}
\label{sec:primary}
\label{sec:ablation}
Adding $g_{12}$ to the 28-D intensity baseline increases cross-track balanced accuracy from $67.03\%$ to $73.33\%$ for RF and from $66.72\%$ to $73.71\%$ for SVC (Table~\ref{tab:benchmark}; Fig.~\ref{fig:primary_contrasts}). The paired increments are $+6.30\pp$ $[5.25,8.19]$ and $+6.99\pp$ $[5.78,9.39]$. Overall accuracy increases by $+1.95\pp$ $[1.32,2.65]$ and $+1.64\pp$ $[0.91,2.47]$. Intervals use 20{,}000 joint block bootstrap draws; the five-class analysis retains 19{,}306 draws containing all classes. Open water occurs in three evaluation blocks and is absent from the remaining draws.

The representation ladder separates the contributions of intensity texture and coherence (Fig.~\ref{fig:coherence_ablation}). Expanding $I_8$ to $I_{28}$ adds $+2.66\pp$ $[1.97,3.34]$ for RF and $+2.73\pp$ $[1.96,3.66]$ for SVC. Adding raw coherence directly to $I_8$ yields $+8.60\pp$ $[7.24,10.36]$ and $+10.10\pp$ $[8.14,12.44]$. The twenty added intensity features therefore supply $31\%$ and $27\%$ of the increments obtained from the single coherence scalar. The compact $I_8{+}g_{12}$ representation reaches $72.97\%$ and $74.09\%$ balanced accuracy, respectively, compared with $73.33\%$ and $73.71\%$ for $I_{28}{+}g_{12}$.

The three derived transformations provide no additional gain. Adding $\sigma_{3\times3}(g_{12})$, $\|\nabla g_{12}\|$, and $\ln(1-g_{12}+10^{-4})$ to $I_8{+}g_{12}$ produces $F_{12}$, whose paired increment is $-0.28\pp$ $[-0.65,-0.08]$ for RF and $-0.20\pp$ $[-0.38,+0.07]$ for SVC. Raw coherence supplies the measured benefit of the interferometric representation.

\begin{figure*}[!t]
\centering
\includegraphics[width=0.98\textwidth]{fig4_primary_contrasts.pdf}
\caption{Cross-track balanced accuracy for five representations: (a) RF and (b) SVC. Markers give full-sample estimates; error bars show one standard deviation across three training-sample seeds. Arrows mark the paired coherence increment over $I_{28}$, with 95\% spatial block-bootstrap intervals.}
\label{fig:primary_contrasts}
\end{figure*}

\begin{table*}[!t]
\caption{Cross-track performance (mean $\pm$ seed standard deviation) and in-track minus cross-track gap $G$. The representations use the same two acquisitions and 400\,m spatial exclusion.}
\label{tab:benchmark}
\centering
\scriptsize
\setlength{\tabcolsep}{3pt}
\begin{tabular}{llcccc}
\toprule
Learner & Repr. & Cross-Track OA (\%) & Cross-Track BA (\%) & $G_{\text{OA}}$ (pp) & $G_{\text{BA}}$ (pp) \\
\midrule
RF & $I_8$ (8-D) & 67.83\,$\pm$\,0.12 & 64.36\,$\pm$\,0.58 & 3.81 & 4.70 \\
RF & $I_{28}$ (28-D) & 69.42\,$\pm$\,0.08 & 67.03\,$\pm$\,0.40 & 3.62 & 4.76 \\
RF & $I_8{+}g_{12}$ (9-D) & 70.29\,$\pm$\,0.08 & 72.97\,$\pm$\,0.42 & 3.48 & 4.35 \\
RF & $F_{12}$ (12-D) & 70.66\,$\pm$\,0.13 & 72.69\,$\pm$\,0.72 & 3.42 & 4.47 \\
RF & $I_{28}{+}g_{12}$ (29-D) & \textbf{71.37}\,$\pm$\,0.06 & \textbf{73.33}\,$\pm$\,0.55 & 3.35 & 4.12 \\
\midrule
SVC & $I_8$ (8-D) & 70.30\,$\pm$\,0.09 & 63.99\,$\pm$\,0.52 & 3.06 & 4.91 \\
SVC & $I_{28}$ (28-D) & 71.13\,$\pm$\,0.05 & 66.72\,$\pm$\,0.14 & 3.31 & 4.74 \\
SVC & $I_8{+}g_{12}$ (9-D) & 72.24\,$\pm$\,0.06 & \textbf{74.09}\,$\pm$\,0.48 & 2.78 & 3.86 \\
SVC & $F_{12}$ (12-D) & 72.23\,$\pm$\,0.03 & 73.89\,$\pm$\,0.28 & 2.83 & 3.89 \\
SVC & $I_{28}{+}g_{12}$ (29-D) & \textbf{72.78}\,$\pm$\,0.03 & 73.71\,$\pm$\,0.43 & 3.12 & 4.44 \\
\bottomrule
\end{tabular}
\\[3pt]
\parbox{\textwidth}{\footnotesize \textbf{Incremental coherence gain} ($I_{28}+g_{12}-I_{28}$), 20{,}000 joint block bootstrap replicates restricted to those in which all five classes are represented (19{,}306 of 20{,}000): RF $\Delta\text{BA} = +6.30\ [5.25, 8.19]\pp$, $\Delta\text{OA} = +1.95\ [1.32, 2.65]\pp$; SVC $\Delta\text{BA} = +6.99\ [5.78, 9.39]\pp$, $\Delta\text{OA} = +1.64\ [0.91, 2.47]\pp$.}
\end{table*}


\begin{figure}[!t]
\centering
\includegraphics[width=0.48\textwidth]{fig14_coherence_ablation.pdf}
\caption{Nested representation contrasts in cross-track balanced accuracy with 95\% joint block bootstrap intervals. The comparisons separate added intensity features, raw coherence, and three derived coherence transformations.}
\label{fig:coherence_ablation}
\end{figure}

\subsection{Transfer Across Tracks}
\label{sec:matrix}
\label{sec:paired}
\label{sec:transfergap}
Coherence improves all twelve directional cross-track transfers for both learners (Supplementary Table~\ref{supp-tab:transfer_matrix}; Fig.~\ref{fig:transfer_matrix}). RF increments range from $+4.7$ to $+7.4\pp$, and SVC increments range from $+6.0$ to $+8.2\pp$. Baseline accuracy depends strongly on the ordered track pair: RF ranges from $60.9\%$ for T139$\to$T37 to $70.8\%$ for T37$\to$T15, while SVC ranges from $59.7\%$ to $70.9\%$. The largest RF increment occurs for T139$\to$T37.



\begin{figure*}[!t]
\centering
\includegraphics[width=0.94\textwidth]{fig5_transfer_matrix.pdf}
\caption{Balanced accuracy for all source--target track pairs under the 400\,m-buffered protocol: (a) RF and (b) SVC. Cells report baseline and augmented accuracy and their paired difference. Dashed outlines identify in-track evaluations.}
\label{fig:transfer_matrix}
\end{figure*}

In-track gains are close to cross-track gains: $+5.66$ versus $+6.30\pp$ for RF and $+6.69$ versus $+6.99\pp$ for SVC. The paired difference $\Delta_{\mathrm{cross}}-\Delta_{\mathrm{in}}$ is $+0.64\pp$ $[+0.07,+1.16]$ in balanced accuracy and $+0.26\pp$ $[-0.05,+0.54]$ in overall accuracy for RF. For SVC, the corresponding values are $+0.30\pp$ $[-0.01,+0.59]$ and $+0.19\pp$ $[-0.17,+0.62]$. Most of the coherence increment is shared by in-track and cross-track classification.

The directional transfer gap decreases from $4.76$ to $4.12\pp$ in RF balanced accuracy and from $4.74$ to $4.44\pp$ in SVC balanced accuracy (Supplementary Fig.~\ref{supp-fig:transfer_gap}). Overall-accuracy gaps decrease from $3.62$ to $3.35\pp$ and from $3.31$ to $3.12\pp$. These reductions equal the paired differences above; only the RF balanced-accuracy interval excludes zero.

Across distinct tracks, incidence-angle separation on $\Omega$ ranges from $3.62^\circ$ for T88$\leftrightarrow$T139 to $12.21^\circ$ for T15$\leftrightarrow$T139. The per-pair coherence increment has no monotone relationship with this separation. Regression summaries over all sixteen pairs, including the in-track anchors, give baseline and augmented slopes of $-0.41$ and $-0.34\pp/^\circ$ for RF and $-0.39$ and $-0.39\pp/^\circ$ for SVC (Supplementary Fig.~\ref{supp-fig:angular_disparity}).





\subsection{Class-Level Discrimination and Spatial Structure}
\label{sec:perclass}
\label{sec:confusion}
Built-up recall accounts for $74.5\%$ of the RF balanced-accuracy increment and $75.5\%$ of the SVC increment (Fig.~\ref{fig:class_decomposition}). Recall rises from $67.22\%$ to $90.71\%$ for RF and from $66.63\%$ to $93.01\%$ for SVC. Forest recall rises from $76.15\%$ to $81.13\%$ and from $78.30\%$ to $82.92\%$. Grassland changes by $+0.48$ and $-0.04\pp$, and cropland by $+1.24$ and $+1.18\pp$.

Equal class weighting assigns built-up gains contributions of $+4.70\pp$ and $+5.28\pp$ to balanced accuracy. Built-up occupies $1.72\%$ of the evaluation sample (165 unique pixels), so its RF contribution to overall accuracy is $+0.40\pp$. This sample-prevalence weighting explains the difference between the balanced-accuracy and overall-accuracy increments.

\begin{figure*}[!t]
\centering
\includegraphics[width=0.94\textwidth]{fig6_class_decomposition.pdf}
\caption{Per-class cross-track recall for RF and SVC, pooled over twelve track pairs and three seeds. Orange circles denote $I_{28}$; blue squares denote $I_{28}+g_{12}$. Connectors show recall changes, with unique pixel counts below class names. Built-up and forest contribute the largest gains.}
\label{fig:class_decomposition}
\end{figure*}

The class signatures show a strong coherence contrast between built-up and vegetation on consensus support $C_2$ (Fig.~\ref{fig:coherence_distributions}). Across tracks, built-up median $g_{12}$ is $0.83$--$0.88$, with median two-date mean VV between $-7.8$--$-7.0$\,dB. Forest has median coherence $0.25$--$0.28$ and median VV between $-9.6$--$-8.7$\,dB. Grassland has median coherence $0.28$--$0.36$ and median VV between $-13.4$--$-12.4$\,dB. Open-water median VV is at most $-19.7$\,dB on every track. The built-up--forest median coherence separation ranges from $0.58$ to $0.62$ across the four viewing configurations.

\begin{figure*}[!t]
\centering
\includegraphics[width=0.98\textwidth]{fig7_coherence_distributions.pdf}
\caption{Class-specific 12-day coherence on consensus support $C_2$: (a) T15 distributions and (b) class medians across the four tracks. Violin densities use at most 10,000 pixels per class; medians and quartiles use all available class pixels. Error bars in (b) show half the interquartile range on either side of the median. On T15, built-up has median coherence $0.83$ and interquartile range $[0.76,0.88]$; its median separation from forest is $0.58$. Track incidence angles span approximately $32.8^\circ$--$45.0^\circ$.}
\label{fig:coherence_distributions}
\end{figure*}

The largest confusion change is the built-up/forest exchange (Table~\ref{tab:confusion}). RF assigns $17.5\%$ of built-up pixels to forest under $I_{28}$ and $0.8\%$ after adding coherence. The reciprocal error decreases from $6.9\%$ to $1.3\%$. SVC shows the same confusion pattern. Cropland remains strongly confused with grassland: the baseline assigns $45.1\%$ of cropland pixels to grassland, and cropland recall changes from $46.0\%$ to $47.2\%$ after augmentation.

\begin{table}[!t]
\caption{Pooled Cross-Track Confusion Matrix for the Random Forest Under the Buffered Protocol, Row-Normalised (\%).
Each Cell Reads $I_{28} \rightarrow I_{28}+g_{12}$; Cells Changing by Less Than $1\pp$ Are Shown Once.}
\label{tab:confusion}
\centering
\scriptsize
\setlength{\tabcolsep}{2.5pt}
\begin{tabular}{lccccc}
\toprule
True $\backslash$ Pred. & Forest & Grass. & Crop. & Urban & Water \\
\midrule
Forest & \textbf{76.2$\to$81.1} & 13.2 & 3.8 & 6.9$\to$1.3 & 0.0 \\
Grassland & 5.3 & \textbf{76.5} & 16.7 & 1.5 & 0.0 \\
Cropland & 5.8 & 45.1 & \textbf{46.0$\to$47.2} & 3.0$\to$1.5 & 0.2 \\
Urban & 17.5$\to$0.8 & 7.7$\to$3.4 & 7.6$\to$5.1 & \textbf{67.2$\to$90.7} & 0.0 \\
Water & 0.7 & 10.1$\to$11.2 & 19.6$\to$16.9 & 0.3 & \textbf{69.3$\to$70.6} \\
\midrule
Support & 69{,}660 & 190{,}872 & 77{,}040 & 5{,}940 & 2{,}088 \\
\bottomrule
\end{tabular}
\\[2pt]
\parbox{\columnwidth}{\footnotesize Support is the cumulative number of predictions over 12 ordered track pairs and 3 seeds ($36\times$ the underlying unique evaluation pixels: 1{,}935 / 5{,}302 / 2{,}140 / 165 / 58). The block bootstrap resamples entire geographic blocks with replacement, carrying all 36 repetitions together, so this pooled support does not inflate the effective sample size.}
\end{table}

Spatial maps show how these class-level changes alter settlement delineation. In the illustrated Wierden/Almelo region spanning evaluation blocks 205 and 206, coherence corrects 176 of 192 false forest predictions over settlement pixels, a $91.7\%$ reduction (Fig.~\ref{fig:spatial_maps}). In the suburban fringe of Block~803, the illustrated T15$\to$T37 transfer gains $24.2\pp$ in local built-up recall (Supplementary Fig.~\ref{supp-fig:multilandscape_maps}). The rural landscape in Block~603 gains $37.5\pp$, with farmsteads and infrastructure recovered in the augmented map. These examples display the spatial pattern of the pooled confusion changes.

\begin{figure*}[!t]
\centering
\includegraphics[width=0.98\textwidth]{fig8_spatial_classification.pdf}
\caption{Spatial classification maps and pixel-level error resolution over real Sentinel-1 SAR observations in evaluation blocks 205 and 206 (Wierden/Almelo, Overijssel, Netherlands; $1.8 \times 2.6$\,km macro domain, $1.2 \times 1.2$\,km inset ROI). (a) False-color SAR observation ($\gamma^0_{\mathrm{VV}}$, $\gamma^0_{\mathrm{VH}}$, $\gamma_{12}$). (b) Harmonized consensus ground truth ($C_2$). (c) Baseline $I_{28}$ predictions exhibiting widespread false forest confusion over urban structures. (d) Coherence-augmented $I_{28}{+}g_{12}$ predictions restoring urban settlements. (e--h) Zoomed ROI details showing: (e) high-resolution SAR composite, (f) 12-day InSAR coherence heatmap ($\gamma_{12}$), (g) baseline error map highlighting 192 false forest pixels, and (h) coherence recovery map correcting 176 pixels ($91.7\%$ confusion reduction).}
\label{fig:spatial_maps}
\end{figure*}



\label{sec:labelaudit}
The 165 built-up evaluation pixels occupy 16 connected objects in eight of the twelve evaluation blocks. One block contains $38.8\%$ of these pixels; the largest object contains $53.3\%$, and the five largest contain $88.5\%$ (Supplementary Fig.~\ref{supp-fig:labelaudit}). At 40\,m posting, $55.8\%$ lie within 40\,m of a class boundary and $78.8\%$ within 80\,m. Fully built-up $3\times3$ and $5\times5$ neighbourhoods occur for $12.7\%$ and $3.0\%$ of the sampled pixels. The minimum-mapping-unit filter removes $16.4\%$ of built-up label pixels, compared with $1.4\%$ of grassland pixels.



\subsection{Robustness to Support and Reference Labels}
\label{sec:label_independence}
North, south, and east directional holdouts retain positive coherence increments with a 5\,km excluded buffer. RF balanced-accuracy increments are $5.83$, $7.75$, and $7.27\pp$, respectively; SVC increments are $8.52$, $8.46$, and $8.07\pp$. These separate spatial partitions extend the comparison to contiguous held-out regions.

Coherence gains remain positive on interior built-up pixels and across object exclusions (Table~\ref{tab:builtup_audit}). Restricting evaluation to the 50 pixels at least 80\,m from a class boundary gives built-up recall increments of $+15.33\pp$ for RF and $+17.44\pp$ for SVC. On the 21 pixels with a fully built-up $3\times3$ neighbourhood, the increments are $+11.90$ and $+15.08\pp$. Both restrictions reduce the increment relative to the full built-up sample.

Resampling the 16 connected objects while holding the other class recalls fixed gives five-class increments of $+6.30\pp$ $[5.56,8.32]$ and $+6.99\pp$ $[6.01,9.70]$. Leaving each object out in turn produces ranges of $+6.02$ to $+7.42\pp$ for RF and $+6.61$ to $+8.46\pp$ for SVC. Removing the largest sampled object increases the increments to $+7.42$ and $+8.46\pp$.

Excluding built-up from the metric gives four-class balanced-accuracy increments of $+2.01\pp$ $[1.50,2.63]$ for RF and $+2.14\pp$ $[1.62,2.73]$ for SVC, with forest providing the largest remaining recall gain. At block level, the five-class increment is positive in all twelve blocks for both learners (two-sided sign test, $p=0.0005$). RF block increments range from $+0.97$ to $+14.95\pp$, with mean $+5.81\pp$; SVC increments range from $+1.07$ to $+17.29\pp$, with mean $+6.27\pp$ (Supplementary Fig.~\ref{supp-fig:spatial_sign_test}).

\begin{table}[!t]
\caption{Audit of the 165 Built-Up Evaluation Pixels. Recall Values Are Cross-Track, Pooled Over
Twelve Ordered Track Pairs and Three Seeds.}
\label{tab:labelaudit}\label{tab:builtup_audit}
\centering
\scriptsize
\setlength{\tabcolsep}{3pt}
\begin{tabular}{lcccc}
\toprule
& \multicolumn{2}{c}{Random forest} & \multicolumn{2}{c}{SVC} \\
\cmidrule(lr){2-3}\cmidrule(lr){4-5}
Subset / test & $I_{28}$ & $I_{28}{+}g_{12}$ & $I_{28}$ & $I_{28}{+}g_{12}$ \\
\midrule
\multicolumn{5}{l}{\emph{Built-up recall (\%)}} \\
All 165 pixels                          & 67.22 & \textbf{90.71} & 66.63 & \textbf{93.01} \\
$\geq 80$\,m from a class boundary (50) & 80.33 & 95.67 & 79.11 & 96.56 \\
Fully built-up $3\times3$ (21)          & 84.79 & 96.69 & 84.13 & 99.21 \\
\midrule
\multicolumn{5}{l}{\emph{Five-class $\Delta\text{BA}$ (pp) under alternative treatments}} \\
Block bootstrap    & \multicolumn{2}{c}{$+6.30$ $[5.25, 8.19]$} & \multicolumn{2}{c}{$+6.99$ $[5.78, 9.39]$} \\
Object-clustered bootstrap   & \multicolumn{2}{c}{$+6.30$ $[5.56, 8.32]$} & \multicolumn{2}{c}{$+6.99$ $[6.01, 9.70]$} \\
Leave-one-object-out range   & \multicolumn{2}{c}{$+6.02$ to $+7.42$}     & \multicolumn{2}{c}{$+6.61$ to $+8.46$} \\
Largest object (88 px) removed & \multicolumn{2}{c}{$+7.42$}              & \multicolumn{2}{c}{$+8.46$} \\
Built-up class removed (4-class BA) & \multicolumn{2}{c}{$+2.01$ $[1.50, 2.63]$} & \multicolumn{2}{c}{$+2.14$ $[1.62, 2.73]$} \\
12-block sign test ($k/12 > 0$)     & \multicolumn{2}{c}{12/12 ($p = 0.0005$)} & \multicolumn{2}{c}{12/12 ($p = 0.0005$)} \\
\bottomrule
\end{tabular}
\\[2pt]
\parbox{\columnwidth}{\footnotesize The object-clustered bootstrap resamples the 16 connected built-up objects with replacement and varies only the built-up recall term, the other four class recalls being held at their point estimates; the block bootstrap varies all five class recalls. For the 4-class balanced accuracy (built-up excluded), 95\% confidence intervals are derived from 20{,}000 paired block bootstrap draws conditioned on complete 4-class presence (19{,}306 valid replicates, 3.47\% whole-class water dropout excluded, identical to the 5-class protocol).}
\end{table}



The label-source analysis also gives positive increments for every training/evaluation label combination (Table~\ref{tab:label_independence_controls}). Under the buffered label-sensitivity protocol, WC$\to$WC yields $+3.64\pp$ for RF and $+6.68\pp$ for SVC on common valid support $U_2$. CLC$\to$CLC yields $+6.53$ and $+8.46\pp$. The cross-product evaluations give $+5.54$ and $+8.76\pp$ for WC$\to$CLC, and $+4.14$ and $+5.69\pp$ for CLC$\to$WC.

On consensus support $C_2$, consensus-trained models yield $+7.02$ and $+7.60\pp$, while WC-trained and WC-evaluated models yield $+6.96$ and $+11.07\pp$. The analysis uses $U_2=600{,}317$ valid pixels on $\Omega=608{,}825$ and $C_2=287{,}013$ consensus pixels. This analysis samples up to 800 evaluation pixels per block on $U_2$, then restricts that sample to $C_2$ for consensus evaluations. Each model uses 20{,}000 training pixels, a 400\,m exclusion from the $U_2$ evaluation sample, and the selected R1 hyperparameters. Its consensus evaluation support therefore differs from the primary benchmark. Positive increments are observed with both reference products and their cross-product evaluations.

\begin{table}[!t]
\caption{Cross-Track Balanced Accuracy Increment ($\Delta\text{BA}$, pp) Across Reference Label Products and Support Sets With 400\,m Training--Evaluation Exclusion.}
\label{tab:label_independence_controls}
\centering
\scriptsize
\setlength{\tabcolsep}{1.8pt}
\begin{tabular}{lllcc}
\toprule
Train Label & Eval Label & Support & RF $\Delta\text{BA}$ & SVC $\Delta\text{BA}$ \\
\midrule
WorldCover (WC) & WorldCover (WC) & Common ($U_2$) & $+3.64$ & $+6.68$ \\
WorldCover (WC) & CORINE (CLC)    & Common ($U_2$) & $+5.54$ & $+8.76$ \\
CORINE (CLC)    & WorldCover (WC) & Common ($U_2$) & $+4.14$ & $+5.69$ \\
CORINE (CLC)    & CORINE (CLC)    & Common ($U_2$) & $+6.53$ & $+8.46$ \\
Consensus ($C_2$) & Consensus ($C_2$) & Consensus ($C_2$) & $+7.02$ & $+7.60$ \\
WorldCover (WC) & WorldCover (WC) & Consensus ($C_2$) & $+6.96$ & $+11.07$ \\
\bottomrule
\end{tabular}
\\[2pt]
\parbox{\columnwidth}{\footnotesize Mean paired increments across twelve ordered cross-track pairs and three seeds. Models use frozen R1 hyperparameters and 20{,}000 training pixels. Consensus evaluations use the agreeing subset of the sampled $U_2$ evaluation pixels.}
\end{table}

\subsection{Temporal Replication}
\label{sec:temporal_validation}
The coherence increment remains positive in balanced accuracy for the October 2024 acquisition window (Table~\ref{tab:temporal_validation}). Across the same four tracks and twelve evaluation blocks, RF balanced accuracy increases from $62.91\%$ to $67.93\%$ ($+5.03\pp$), and SVC increases from $63.01\%$ to $68.24\%$ ($+5.23\pp$). The corresponding 95\% bias-corrected and accelerated block bootstrap intervals are $[3.33,7.08]$ and $[3.36,7.27]\pp$.

Autumn overall-accuracy increments are $+0.62\pp$ for RF and $-0.01\pp$ for SVC, with percentile intervals $[-0.36,1.74]$ and $[-0.68,0.79]\pp$. Built-up recall increases from $75.00\%$ to $96.08\%$ for RF and from $73.82\%$ to $97.04\%$ for SVC, giving gains of $+21.08$ and $+23.22\pp$. Forest recall gains are $+4.10$ and $+1.05\pp$.

\begin{table*}[!t]
\caption{Temporal Replication Benchmark: Spring 2024 vs. Autumn 2024 Cross-Track Evaluation Under Observation Budget Parity ($N=2$)}
\label{tab:temporal_validation}
\centering
\scriptsize
\setlength{\tabcolsep}{6pt}
\begin{tabular}{lcccc}
\toprule
& \multicolumn{2}{c}{Spring 2024} & \multicolumn{2}{c}{Autumn 2024 (Temporal Replication)} \\
\cmidrule(lr){2-3}\cmidrule(lr){4-5}
Evaluation Metric / Diagnostic Quantity & Random Forest & SVC & Random Forest & SVC \\
\midrule
Five-Class Balanced Accuracy Increment ($\Delta\text{BA}$, pp) & $+6.30$ & $+6.99$ & $+5.03$ & $+5.23$ \\
\quad 95\% Joint Block Bootstrap Interval                      & $[5.25, 8.19]$ & $[5.78, 9.39]$ & $[3.33, 7.08]$ & $[3.36, 7.27]$ \\
Sample Overall Accuracy Increment ($\Delta\text{OA}$, pp)& $+1.95$ & $+1.64$ & $+0.62$ & $-0.01$ \\
\quad 95\% Joint Block Bootstrap Interval                      & $[1.32, 2.65]$ & $[0.91, 2.47]$ & $[-0.36, 1.74]$ & $[-0.68, 0.79]$ \\
Built-Up Class Recall Gain ($\Delta\text{Recall}$, pp)           & $+23.48$ & $+26.38$ & $+21.08$ & $+23.22$ \\
\quad Absolute Recall Shift (Baseline $\to$ Augmented)         & $67.2\% \to 90.7\%$ & $66.6\% \to 93.0\%$ & $75.0\% \to 96.1\%$ & $73.8\% \to 97.0\%$ \\
Four-Class Balanced Accuracy Increment (Without Built-Up, pp)             & $+2.01$ & $+2.14$ & $+1.01$ & $+0.74$ \\
12-Block Spatial Consistency (Sign Test $k/12 > 0$)             & 12/12 ($p = 0.0005$) & 12/12 ($p = 0.0005$) & 9/12 ($p = 0.146$) & 10/12 ($p = 0.0386$) \\
\bottomrule
\end{tabular}
\\[2pt]
\parbox{\textwidth}{\footnotesize All metrics evaluated across all 12 ordered cross-track pairs ($n=12$) on the identical 12 evaluation blocks and $C_2$ consensus ground-truth support under strict observation budget parity ($N=2$). Spring intervals use the joint block percentile bootstrap. Autumn balanced-accuracy intervals use the bias-corrected and accelerated bootstrap; autumn overall-accuracy intervals use the percentile bootstrap (20{,}000 draws).}
\end{table*}

Removing built-up leaves four-class balanced-accuracy increments of $+1.01\pp$ for RF and $+0.74\pp$ for SVC, compared with $+2.01$ and $+2.14\pp$ in spring. Positive block increments occur in 9 of 12 autumn blocks for RF ($p=0.146$) and 10 of 12 for SVC ($p=0.0386$). The autumn results retain the concentration of gains in built-up classification, with smaller five-class increments and weaker spatial consistency than in spring.
